\section{Die These}

\textbf{[SETZUNG]} Das Standardmodell ist in dieser Theorie kein konkurrierender
Rahmen, sondern die \emph{dekompaktifizierte Projektion} der kompakten Struktur.

Was im Standardmodell als Liste unabhängiger Eingaben erscheint -- Massen,
Kopplungen, Mischungswinkel --, sind in dieser Lesart Projektionsbilder einer
Struktur mit weniger Freiheitsgraden.

\section{Was die Projektion erklärt}

\textbf{[B]} Drei Züge des Standardmodells werden in dieser Lesart zu Folgen
und nicht zu Annahmen:

\textbf{Die Zahl drei.} Drei Generationen folgen aus der $Z_3$-Struktur
(A030, A110). Im Standardmodell ist sie eine Beobachtung ohne Grund.

\textbf{Die Hierarchie.} Die Größenordnungen der Leptonmassen folgen der
$\xi$-Leiter (A100), auf Prozentniveau. Im Standardmodell sind die
Yukawa-Kopplungen freie Parameter ohne Ordnung untereinander.

\textbf{Die Form der Massenrelation.} Die Koide-Relation $Q=2/3$ ist
Operatorbedingung, nicht Zufall (A110). Im Standardmodell hat sie keine Stellung.

\section{Was die Projektion nicht erklärt}

\textbf{[SETZUNG]} Dieser Abschnitt ist der längere, und das ist der ehrliche
Befund.

\textbf{Die Eichgruppe — drei konvergente Argumente.}

\textbf{[B] Argument 1: Ladungsquantisierung aus Torus-Fluss.}
Die elektrische Ladung entsteht aus topologisch quantisierten Flusslinien
durch den Torus: $\Phi=n\cdot h/e$ mit $n\in\mathbb{Z}$. Die Quantisierung ist
topologisch geschützt -- der Torus hat zwei nicht-kontrahierbare Schleifen
(toroidal und poloidal), der Fluss durch sie ist invariant unter stetigen
Deformationen. $U(1)_\text{EM}$ folgt damit aus der Torus-Topologie, nicht
aus einer postulierten Eichsymmetrie.

\textbf{[B] Argument 2: Farbladung aus Verschlingungszahl.}
Die Farbladung der QCD entsteht aus der topologischen Verschlingung
(linking number) von genau drei Flussfäden um den Torus. Drei Farben
(rot, grün, blau) = drei Fäden: $SU(3)_C$ ist topologisch erzwungen
durch die Dreifach-Verschlingungsstruktur des $T^4/\mathbb{Z}_3$.

\textbf{[S] Argument 3: SU(27)-Gruppenkettenzerlegung.}
Aus $27=3^3$ ($\mathbb{Z}_3$-Torus) folgt die Kette
$SU(27)\supset SU(9)^3\supset SU(3)^9\supset SU(3)\times SU(2)\times U(1)$.
Das macht die SM-Eichgruppe zur kanonischen Untergruppe der Symmetriebrechung.

\textbf{Was noch fehlt:} Die Argumente 1 und 2 liefern $U(1)$ und $SU(3)$
aus der Torus-Topologie. $SU(2)_L$ folgt aus der chiralen Projektion (A095).
Was nicht ausgeschrieben ist: eine explizite Abbildung der Torus-Topologie
auf die kovarianten Ableitungen $D_\mu = \partial_\mu + igA_\mu^a T^a$ --
also der Übergang von der topologischen Ladungsstruktur zum Eichlagrangian
$-\tfrac14 F_{\mu\nu}^a F^{a\,\mu\nu}$.

\textbf{Die Kopplungskonstanten.} Ihre Werte und ihr Laufen sind nicht
Gegenstand der A-Serie. Für $\alpha$ ist der Status in A130 geklärt:
kalibriert, nicht hergeleitet.

\textbf{Die Mischungsmatrizen.} CKM und PMNS erscheinen nicht.

\textbf{Der Higgs-Sektor — Konsistenzprüfung nahe Beweis.} Eine vollständige 1-Loop-EFT-Matching-Rechnung (im Altkorpus ausgeführt) leitet ab:
vollständige 1-Loop-EFT-Matching-Rechnung durch und leitet ab:
\[
  \xi = \frac{\lambda_h^2 v^2}{16\pi^3 m_h^2}
\]
mit $\lambda_h=m_h^2/(2v^2)$, $v=246\,\text{GeV}$, $m_h=125{,}09\,\text{GeV}$.
Numerisch: $\xi_\text{EFT}=1{,}303\times10^{-4}$, Abweichung von $\xi_0=4/30000$
nur $2{,}3\,\%$. Die $16\pi^3$-Struktur ist das natürliche Ergebnis der
1-Loop-Unterdrückung und der NDA-Normierung -- kein Fitting.

\textbf{[S]} \textbf{Zwei Einschränkungen, die zusammen entscheidend sind.}

\textbf{Erstens} liefert die EFT-Rechnung ein leptonspezifisches $\xi$
($\xi_e\approx6{,}7\times10^{-14}$, $\xi_\mu\approx2{,}9\times10^{-9}$),
nicht das universelle $\xi_0$. Der numerisch nahe Wert
$\xi_\text{EFT}\approx1{,}303\times10^{-4}$ ist der Wert für ein Referenzfermion
bei $\mu=m_h$ im $\overline{\text{MS}}$-Schema.

\textbf{Zweitens ist die $2{,}3\,\%$-Abweichung strukturell erwartet} -- aus
zwei unabhängigen Gründen, die zusammenwirken:

\textbf{(a) Ausgerollte SI-Werte.} Die Formel $\xi_\text{EFT}=m_h^2/(64\pi^3 v^2)$
enthält $v=246\,\text{GeV}$ und $m_h=125{,}09\,\text{GeV}$ -- SI-Messwerte
im $\overline{\text{MS}}$-Schema, schemaabhängig und nur auf 1-Loop-Niveau.
Eine $2{,}3\,\%$-Abweichung entspricht $1{,}1\,\%$ Abweichung in $v$, was
innerhalb der theoretischen Unsicherheit von 1-Loop-Matchings liegt (typisch
$1$--$5\,\%$ ohne NNLO-Korrekturen).

\textbf{(b) Fraktale Korrekturen sind in $v$ nicht sauber trennbar.}
Der gemessene VEV $v=246\,\text{GeV}$ absorbiert alle SM-Quantenkorrekturen --
die fraktale Torus-Struktur $D_f=3-\xi$ ist darin nicht separat sichtbar,
sondern implizit in der Messung enthalten und mit den SM-Schleifen
vermengt. Das EFT-Matching im flachen Raum mit ganzzahliger Dimension $D=3$
kann $K_\text{frak}=(1-100\xi)$ daher nicht sauber herauslösen, was zu
einer systematischen Ungenauigkeit führt. Rechnerisch zeigt sich das:
$\xi_\text{EFT}/K_\text{frak}^{3/2}\approx1{,}330\times10^{-4}$, Abweichung
von $\xi_0$ nur noch $0{,}3\,\%$ -- aber das ist eine post-hoc-Korrektur,
kein Beweis, da $K_\text{frak}$ in $v$ nicht isoliert vorliegt.

Kurz: die Abweichung ist strukturell erwartet, weil zwei verschiedene
Beschreibungsebenen -- flache perturbative QFT und fraktale Torus-Geometrie
-- nicht denselben effektiven $v$-Wert sehen.

\textbf{$\xi_\text{EFT}$ darf deshalb nicht für die Berechnung von $\alpha$
oder $E_0$ verwendet werden.} Das wäre eine Zirkularitätsfalle: man würde
SI-Messwerte ($v$, $m_h$) durch die FFGFT-Formeln laufen lassen und das
Ergebnis als Bestätigung werten, obwohl die SI-Werte selbst die Abweichung
tragen. Die Formeln der A-Serie verwenden ausschließlich $\xi_0=4/30000$.

$\xi_\text{EFT}$ ist ein Wilson-Koeffizient (scheme- und skalenabhängig);
$\xi_0$ ist das geometrische Verhältnis $L_0/\ell_P$ (scheme-unabhängig).
Dass beide auf $2{,}3\,\%$ übereinstimmen -- bei einer Rechnung die
SI-Messwerte mit eingebauter theoretischer Unsicherheit verwendet --
ist der Konsistenz-Befund. Mehr ist es nicht, aber auch nicht weniger.

\textbf{[B]} Die Bilanz lautet damit: die Theorie erklärt die \emph{Ordnung} des
Fermionsektors und nicht den Eichsektor. Das ist eine Teilaussage, und sie als
Erklärung des Standardmodells auszugeben wäre eine Überdehnung.

\section{Warum die Projektion Information verliert}

\textbf{[B]} Nach A090 ist die räumliche Projektion (Typ~II) nicht umkehrbar.
Daraus folgt, dass das Standardmodell die kompakte Struktur \emph{nicht}
vollständig wiedergeben kann -- unabhängig davon, wie genau man es misst.

Das ist prüfbar, wenn auch schwer: was bei der Projektion nicht vollständig
verschwindet, hinterlässt einen Rest der Ordnung $\xi$, und die vier
Prognosegrößen der Serie sind genau diese Reste (A090, A210).


\section{Prüfskript}

\begin{itemize}[leftmargin=2em,itemsep=2pt,topsep=2pt]
  \item Numerische Verifikation der Projektionsketten-Schritte und
    Massenverhältnisse aus A090/A100:
    \texttt{python/A\_Serie\_Skripte/a090\_proj.py} (gemeinsam mit A090).
\end{itemize}

\section{Was offen bleibt}

\textbf{Erstens fehlt der Eichsektor vollständig.} Ohne ihn ist die Behauptung,
das Standardmodell sei eine Projektion, für den größten Teil des
Standardmodells nicht geprüft.

\textbf{Zweitens ist die Projektionsabbildung nicht ausgeschrieben.} Es gibt in
dieser Serie keine explizite Abbildung von $T^4/Z_3$ auf die Felder des
Standardmodells -- nur die Behauptung, dass es sie gibt, und Beispiele im
Fermionsektor.

\textbf{Drittens ist die Aussage in der vorliegenden Form nicht falsifizierbar.}
Solange die Abbildung fehlt, kann die These weder bestätigt noch widerlegt
werden. Sie ist ein Programm.

\section{Ersetzt}

Dieses Dokument nimmt auf: Dok.~019 (Standardmodell), Dok.~097 (QFT), Dok.~141
(SM-Parameter), Dok.~298 (SM als Projektion).

\section{Verwendung von KI-Werkzeugen}

Dieses Dokument ist unter Verwendung KI-gestützter Werkzeuge entstanden. Die
Fragestellungen, die Setzungen und alle inhaltlichen Entscheidungen stammen vom
Autor; die Werkzeuge formulieren aus, rechnen nach und erstellen Prüfskripte.
Die Verantwortung für Inhalt und Fehler liegt vollständig beim Autor. Der
ausführliche Wortlaut steht in A010.
